% GENERATED FILE. Edit canonical lessons, then run npm run generate:books.
% canonical-source-hash: fnv1a64:6e668f8a404dec94
% canonical-lessons: ZH-C190-shudian, ZH-W190-dian-shop, ZH-C190-yaodian, ZH-C190-shangchang, ZH-C190-dianyingyuan, ZH-C190-bowuguan

\chapter{Shops in Town}
\label{ch:zh-shops-in-town}
\hlchaptermodality{190}

\begin{chapteropening}
\textbf{By the end of this chapter:} \emph{I can name the bookshop, the chemist's, the shopping centre, the cinema and the museum.}

Complete the last lesson of chapter 190: I can name the bookshop, the chemist's, the shopping centre, the cinema and the museum.
\end{chapteropening}

\section[\texorpdfstring{\zh{书店} (shūdiàn)}{书店 (shūdiàn)}]{\zh{书店} (shūdiàn) — bookshop}

\label{lesson:ZH-C190-shudian}



\begin{quote}
Before the new one: say the Mandarin for to sign your name, then the Mandarin for identity card.
\end{quote}

\subsection*{You'll want to know: \zh{书店}}
\textbf{\zh{书店}} — \emph{shūdiàn} — \textquotedblleft{}bookshop\textquotedblright{}.

\begin{quote}
\textbf{\zh{这儿的书店很大}} \emph{zhèr de shūdiàn hěn dà} — The bookshops here are big.
\end{quote}

The new character, in pieces: \zh{店} — eight strokes: a roof of a dot, a bar and a long stroke falling to the left, and under it a vertical with a short stroke to its right, over \zh{口}. The next lesson puts it in your hand.

\begin{cousinweb}
\textbf{\zh{书店}} (\emph{shūdiàn}) is built from \zh{书} 'book' and \zh{店} 'shop', both inherited from Middle Chinese, \emph{syo} and \emph{temH} (Guangyun, in Baxter's transcription). The dictionary explains the second, after an older book, as a place where goods are set out for sale.
\end{cousinweb}

\begin{grammarlens}[title={What you've built}]
One more word for asking where things are.
\end{grammarlens}

\subsection*{Your turn}
\begin{itemize}
\raggedright
  \item \emph{Say it:} \emph{shūdiàn}
  \item \emph{Say it:} \emph{shūdiàn}, once more
  \item \emph{Say it:} \emph{zhèr de shūdiàn hěn dà} — The bookshops here are big.
  \item \emph{Recall:} say the Mandarin for to sign your name, then the Mandarin for identity card, then say \emph{shūdiàn} again
\end{itemize}

\begin{tcolorbox}[breakable,colback=teal!4,colframe=teal!35!black,title={Before you move on}]
What does \zh{书店} mean? (\textquotedblleft{}Bookshop\textquotedblright{}.) Say it once more.
\end{tcolorbox}
\section[\texorpdfstring{\zh{店} (diàn)}{店 (diàn)}]{\zh{店} — a roof over a small piece and \zh{口}}

\label{lesson:ZH-W190-dian-shop}



\begin{quote}
The word first, then the character inside it.

\begin{itemize}
\raggedright
  \item \emph{Recall:} say \emph{shūdiàn} — \textbf{R1}, one lesson back, before it decays
\end{itemize}

\begin{quote}
\zh{书} \zh{口}
\end{quote}

\emph{Read it:} each one aloud, with nothing said first — \emph{shū}, \emph{kǒu}
\end{quote}

\subsection*{Script — a roof over a small piece and \zh{口}}
\hlblockfigure{\detokenize{figures/ZH-W190-dian-shop-filmstrip.pdf}}{How \zh{店} is written, stroke by stroke}

\begin{quote}
\zh{店}
\end{quote}

\textbf{Three pieces: a roof, and under it a small piece over \zh{口}.}

\textbf{The roof, three strokes.} A dot at the top, a bar across, then a long stroke from the left end of the bar, falling down to the left all the way to the foot.

\textbf{The small piece, two strokes.} A vertical under the roof, then a short stroke from its middle out to the right.

\textbf{\zh{口}, three strokes:} a short vertical, a stroke that runs right and turns down, then the bar that closes the foot.

3 + 2 + 3 = eight.

\subsection*{Your turn}
\begin{itemize}
\raggedright
  \item \emph{Write it:} \zh{店}, all eight, with the model in front of you
  \item \emph{Cover:} write it again from memory, and count the pieces, not the strokes
\end{itemize}

\begin{tcolorbox}[breakable,colback=teal!4,colframe=teal!35!black,title={Before you move on}]
How many strokes in \zh{店}? (\textbf{Eight}.) Which stroke falls all the way to the foot? (\textbf{The third, the long one on the left}.) What closes the character? (\textbf{The foot bar of \zh{口}}.)

Source: \href{https://raw.githubusercontent.com/chanind/hanzi-writer-data/68d10a4b21150cae5e1ebbd223eed289cf32d90c/data/\%E5\%BA\%97.json}{Hanzi Writer Data \zh{店}, pinned stroke paths and medians}.
\end{tcolorbox}
\section[yàodiàn]{yàodiàn — chemist's; pharmacy}

\label{lesson:ZH-C190-yaodian}



\begin{quote}
Before the new one: say the Mandarin for identity card, then the Mandarin for bookshop.
\end{quote}

\subsection*{You'll want to know: yàodiàn}
\emph{yàodiàn} — \textquotedblleft{}chemist's; pharmacy\textquotedblright{}.

\begin{quote}
\emph{yàodiàn jǐ diǎn kāimén} — What time does the chemist's open?
\end{quote}

Its characters come later; for now it is a word for the ear.

\begin{cousinweb}
\emph{yàodiàn} is built from \emph{yào} 'medicine' and \zh{店} 'shop', both inherited from Middle Chinese, \emph{yak} and \emph{temH} (Guangyun, in Baxter's transcription); the dictionary glosses the first as a herb that cures illness.
\end{cousinweb}

\begin{grammarlens}[title={What you've built}]
One more word for asking where things are.
\end{grammarlens}

\subsection*{Your turn}
\begin{itemize}
\raggedright
  \item \emph{Say it:} \emph{yàodiàn}
  \item \emph{Say it:} \emph{yàodiàn}, once more
  \item \emph{Say it:} \emph{yàodiàn jǐ diǎn kāimén} — What time does the chemist's open?
  \item \emph{Recall:} say the Mandarin for identity card, then the Mandarin for bookshop, then say \emph{yàodiàn} again
\end{itemize}

\begin{tcolorbox}[breakable,colback=teal!4,colframe=teal!35!black,title={Before you move on}]
What does yàodiàn mean? (\textquotedblleft{}Chemist's; pharmacy\textquotedblright{}.) Say it once more.
\end{tcolorbox}
\section[shāngchǎng]{shāngchǎng — shopping centre}

\label{lesson:ZH-C190-shangchang}



\begin{quote}
Before the new one: say the Mandarin for bookshop, then the Mandarin for chemist's.
\end{quote}

\subsection*{You'll want to know: shāngchǎng}
\emph{shāngchǎng} — \textquotedblleft{}shopping centre\textquotedblright{}.

\begin{quote}
\emph{zhōumò de shāngchǎng hěn rènao} — The shopping centre is busy at the weekend.
\end{quote}

Its characters come later; for now it is a word for the ear.

\begin{cousinweb}
\emph{shāngchǎng} is built from \emph{shāng} 'trade' and \emph{chǎng} 'an open ground', both inherited from Middle Chinese, \emph{syang} and \emph{drjang} (Guangyun, in Baxter's transcription). The dictionary glosses the second as a cleared ground for threshing grain, with a level tone, so the rising tone of chǎng is not the regular outcome; the sense 'trade' for the first is not given there either, so how both came is not settled.
\end{cousinweb}

\begin{grammarlens}[title={What you've built}]
One more word for asking where things are.
\end{grammarlens}

\subsection*{Your turn}
\begin{itemize}
\raggedright
  \item \emph{Say it:} \emph{shāngchǎng}
  \item \emph{Say it:} \emph{shāngchǎng}, once more
  \item \emph{Say it:} \emph{zhōumò de shāngchǎng hěn rènao} — The shopping centre is busy at the weekend.
  \item \emph{Recall:} say the Mandarin for bookshop, then the Mandarin for chemist's, then say \emph{shāngchǎng} again
\end{itemize}

\begin{tcolorbox}[breakable,colback=teal!4,colframe=teal!35!black,title={Before you move on}]
What does shāngchǎng mean? (\textquotedblleft{}Shopping centre\textquotedblright{}.) Say it once more.
\end{tcolorbox}
\section[diànyǐngyuàn]{diànyǐngyuàn — cinema}

\label{lesson:ZH-C190-dianyingyuan}



\begin{quote}
Before the new one: say the Mandarin for chemist's, then the Mandarin for shopping centre.
\end{quote}

\subsection*{You'll want to know: diànyǐngyuàn}
\emph{diànyǐngyuàn} — \textquotedblleft{}cinema\textquotedblright{}.

\begin{quote}
\emph{diànyǐngyuàn zài shāngchǎng lǐ} — The cinema is inside the shopping centre.
\end{quote}

Its characters come later; for now it is a word for the ear.

\begin{cousinweb}
\emph{diànyǐngyuàn} is built from \emph{diànyǐng} 'film' and \emph{yuàn} 'a walled courtyard', the \emph{yuàn} of \emph{yīyuàn}; all three parts are inherited from Middle Chinese, \emph{denH}, \emph{'jaengX} and \emph{hjwenH} (Guangyun, in Baxter's transcription).
\end{cousinweb}

\begin{grammarlens}[title={What you've built}]
One more word for asking where things are.
\end{grammarlens}

\subsection*{Your turn}
\begin{itemize}
\raggedright
  \item \emph{Say it:} \emph{diànyǐngyuàn}
  \item \emph{Say it:} \emph{diànyǐngyuàn}, once more
  \item \emph{Say it:} \emph{diànyǐngyuàn zài shāngchǎng lǐ} — The cinema is inside the shopping centre.
  \item \emph{Recall:} say the Mandarin for chemist's, then the Mandarin for shopping centre, then say \emph{diànyǐngyuàn} again
\end{itemize}

\begin{tcolorbox}[breakable,colback=teal!4,colframe=teal!35!black,title={Before you move on}]
What does diànyǐngyuàn mean? (\textquotedblleft{}Cinema\textquotedblright{}.) Say it once more.
\end{tcolorbox}
\section[bówùguǎn]{bówùguǎn — museum}

\label{lesson:ZH-C190-bowuguan}



\begin{quote}
Before the new one: say the Mandarin for shopping centre, then the Mandarin for cinema.
\end{quote}

\subsection*{You'll want to know: bówùguǎn}
\emph{bówùguǎn} — \textquotedblleft{}museum\textquotedblright{}.

\begin{quote}
\emph{bówùguǎn jīntiān bù kāimén} — The museum is closed today.
\end{quote}

Its characters come later; for now it is a word for the ear.

\begin{cousinweb}
\emph{bówùguǎn} is built from \emph{bó} 'broad', \emph{wù} 'a thing' and \emph{guǎn} 'a lodging, a hall', all inherited from Middle Chinese, \emph{pak}, \emph{mjut} and \emph{kwanH} (Guangyun, in Baxter's transcription): a hall of all kinds of things.
\end{cousinweb}

\begin{grammarlens}[title={What you've built}]
One more word for asking where things are.
\end{grammarlens}

\subsection*{Your turn}
\begin{itemize}
\raggedright
  \item \emph{Say it:} \emph{bówùguǎn}
  \item \emph{Say it:} \emph{bówùguǎn}, once more
  \item \emph{Say it:} \emph{bówùguǎn jīntiān bù kāimén} — The museum is closed today.
  \item \emph{Recall:} say the Mandarin for shopping centre, then the Mandarin for cinema, then say \emph{bówùguǎn} again
\end{itemize}

\begin{tcolorbox}[breakable,colback=teal!4,colframe=teal!35!black,title={Before you move on}]
What does bówùguǎn mean? (\textquotedblleft{}Museum\textquotedblright{}.) Say it once more.
\end{tcolorbox}
